\documentclass[11pt]{article}

\usepackage[margin=1in]{geometry}
\usepackage{amsmath,amssymb,amsthm}
\usepackage{graphicx}
\usepackage[hidelinks]{hyperref}

\newtheorem{theorem}{Theorem}
\newtheorem{lemma}{Lemma}
\newtheorem{remark}{Remark}

\newcommand{\E}{\mathbb{E}}
\newcommand{\R}{\mathbb{R}}
\newcommand{\one}{\mathbf{1}}

\title{A \(0<q\le 5\) partial solution to the non-negative Hunter moment conjecture}
\author{Run \texttt{fa\_banach\_001}, lane 5}
\date{2026-07-18}

\begin{document}
\maketitle

\begin{abstract}
Brazitikos and Pandis conjecture that, for non-negative coefficients
\(\sum_i a_i^2=1\), the real moment
\(\E(a_1X_1+\cdots+a_nX_n)^q\) is bounded below by the least equal-support
moment \(\min_{1\le m\le n}\rho(m,q)\).  This packet proves the conjecture
for all dimensions in the range \(0<q\le 5\).  The new step beyond the known
\(q\le 4\) Schur-concavity range is to combine the source paper's integer
endpoint \(k=5\) theorem with Borell log-concavity of normalized moments.
\end{abstract}

\section{Source question}

The source is:

\begin{quote}
Silouanos Brazitikos and Christos Pandis, \emph{Sharp inequalities for symmetric
polynomials, Hunter's conjecture, and moments of exponential random variables},
arXiv:2512.12254.
\end{quote}

The relevant open problem is Conjecture 4.8 in the source PDF:

\begin{center}
\includegraphics[width=0.92\textwidth]{figures/open_problem_crop.png}
\end{center}

For \(m\in\mathbb{N}\) and \(q>0\), write
\[
 \rho(m,q)=\E\left(\frac{X_1+\cdots+X_m}{\sqrt m}\right)^q
 = \frac{\Gamma(m+q)}{\Gamma(m)m^{q/2}},
\]
where the \(X_i\) are i.i.d. standard exponential random variables.

\section{Proof intuition}

The source paper proves the conjectured lower bound at every non-negative
integer exponent.  At the integer \(k=5\), the least equal-support moment is
still the one-coordinate value \(\rho(1,5)=\Gamma(6)\).  On the other hand, any
positive linear combination of independent exponentials has a log-concave
density, and Borell's lemma says that
\[
q\mapsto \frac{\E Y^q}{\Gamma(q+1)}
\]
is log-concave.  Interpolating between \(q=0\), where the normalized moment is
\(1\), and \(q=5\), where the source's integer theorem gives a value at least
\(1\), forces the normalized moment to stay at least \(1\) throughout
\([0,5]\).  A separate scalar lemma checks that, throughout the same interval,
\(\rho(1,q)\) is exactly the minimum among the \(\rho(m,q)\).

\section{Scalar minimum through \(q=5\)}

\begin{lemma}\label{lem:rho-min}
For every integer \(m\ge 1\) and every \(0\le q\le 5\),
\[
\rho(m,q)\ge \rho(1,q)=\Gamma(q+1).
\]
Consequently, for every \(n\ge 1\),
\[
\min_{1\le m\le n}\rho(m,q)=\rho(1,q),\qquad 0\le q\le 5.
\]
\end{lemma}

\begin{proof}
For \(m=1\) there is equality.  Fix \(m\ge 2\) and set
\[
R_m(q)=\frac{\rho(m,q)}{\Gamma(q+1)}
 =\frac{\Gamma(m+q)}{\Gamma(m)\Gamma(q+1)m^{q/2}}.
\]
The function \(\log R_m(q)\) is concave on \([0,\infty)\), because
\[
\frac{d^2}{dq^2}\log R_m(q)=\psi_1(m+q)-\psi_1(q+1)\le 0,
\]
where \(\psi_1\) is the trigamma function, which is decreasing on
\((0,\infty)\).  Since \(R_m(0)=1\), it remains enough to show \(R_m(5)\ge 1\).

At \(q=5\),
\[
R_m(5)=\frac{(m+1)(m+2)(m+3)(m+4)}{120m^{3/2}}.
\]
For real \(m\ge 2\), let the right-hand side be \(r(m)\).  Then
\[
\frac{d}{dm}\log r(m)
=\sum_{j=1}^4\frac{1}{m+j}-\frac{3}{2m}
=\frac{5m^4+30m^3+35m^2-50m-72}
{2m(m+1)(m+2)(m+3)(m+4)} >0
\]
for \(m\ge 2\).  Thus \(r(m)\ge r(2)=3/(2\sqrt2)>1\) for all \(m\ge 2\).
By concavity of \(\log R_m\),
\[
\log R_m(q)\ge \left(1-\frac q5\right)\log R_m(0)
 + \frac q5\log R_m(5)\ge 0
\]
for \(0\le q\le 5\).  This proves the lemma.
\end{proof}

\section{Partial solution}

\begin{theorem}\label{thm:q-le-five}
Let \(0<q\le 5\), let \(a_1,\ldots,a_n\ge 0\) satisfy
\(\sum_{i=1}^n a_i^2=1\), and let \(X_1,\ldots,X_n\) be i.i.d. standard
exponential random variables.  Then
\[
\E(a_1X_1+\cdots+a_nX_n)^q
\ge \min_{1\le m\le n}\rho(m,q).
\]
In fact, the stronger bound
\[
\E(a_1X_1+\cdots+a_nX_n)^q\ge \Gamma(q+1)
\]
holds throughout this range.
\end{theorem}

\begin{proof}
Discard zero coefficients and put
\[
Y=a_1X_1+\cdots+a_nX_n.
\]
For \(a_i>0\), the random variable \(a_iX_i\) has log-concave density
\[
t\mapsto a_i^{-1}e^{-t/a_i}\one_{[0,\infty)}(t).
\]
By the Prekopa-Leindler theorem, convolution preserves log-concavity, so \(Y\)
has a log-concave density on \((0,\infty)\).

The source paper's Lemma 4.13, a Borell log-concavity lemma, therefore applies:
\[
G_Y(s):=\frac{\E Y^s}{\Gamma(s+1)}
\]
is log-concave on \((-1,\infty)\).  We have \(G_Y(0)=1\).

By Theorem 4.5 of the source paper, applied at the integer exponent \(k=5\),
\[
\E Y^5\ge \min_{1\le m\le n}\rho(m,5).
\]
Lemma~\ref{lem:rho-min} gives \(\min_m\rho(m,5)=\rho(1,5)=\Gamma(6)\), hence
\[
G_Y(5)=\frac{\E Y^5}{\Gamma(6)}\ge 1.
\]
For \(0\le q\le 5\), log-concavity now yields
\[
G_Y(q)\ge G_Y(0)^{1-q/5}G_Y(5)^{q/5}\ge 1.
\]
Thus \(\E Y^q\ge \Gamma(q+1)\).  A final application of
Lemma~\ref{lem:rho-min} identifies \(\Gamma(q+1)=\rho(1,q)\) as
\(\min_{1\le m\le n}\rho(m,q)\), proving the conjectured inequality in the
claimed range.
\end{proof}

\section{Verification and novelty notes}

The auxiliary script \texttt{code/check\_q\_le\_5\_endpoint.py} numerically
checks the scalar endpoint ratios \(R_m(5)\), verifies the finite-grid
minimum \(\rho(m,q)/\Gamma(q+1)\ge 1\) for \(0\le q\le 5\), and records that
the naive integer-endpoint interpolation strategy fails just beyond \(q=5\).
These computations are not used as proof.

Local index search found no previous packet for arXiv:2512.12254 or this
specific \(q\le 5\) consequence.  A bounded web/arXiv search on 2026-07-18
found the later adjacent preprint arXiv:2602.03058, \emph{Moments of sums of
exponentials, beyond CHS}.  Its Theorem 2 proves Schur-concavity of the
non-negative moment map for \(0\le p\le 4\), which gives the source conjecture
in the already-known range \(q\le 4\).  The present packet extends the
all-dimensional verified range to \(q\le 5\) by a different endpoint argument.

\section{Limitations}

This is not a full solution of Conjecture 4.8.  The interval \(q>5\) remains
open here.  The same Borell interpolation method cannot simply continue from
the integer endpoints \(5\) and \(6\): numerically, near \(q=5.35\), the chord
between the known integer endpoint lower bounds is below the conjectured
minimum of the \(\rho(m,q)\) curve.

\section{Human review recommendation}

This packet should be reviewed as a likely valid partial solution.  The most
important checks are the invocation of Theorem 4.5 at \(k=5\), the
log-concavity of the convolution density for positive coefficients, and the
scalar Lemma~\ref{lem:rho-min}.

\begin{thebibliography}{9}

\bibitem{BP}
S. Brazitikos and C. Pandis,
\emph{Sharp inequalities for symmetric polynomials, Hunter's conjecture, and
moments of exponential random variables},
arXiv:2512.12254, 2025.

\bibitem{BTT}
S. Brazitikos, C. Tang, and T. Tkocz,
\emph{Moments of sums of exponentials, beyond CHS},
arXiv:2602.03058, 2026.

\bibitem{Prekopa}
A. Prekopa,
Logarithmic concave measures with application to stochastic programming,
\emph{Acta Scientiarum Mathematicarum} 32 (1971), 301--316.

\end{thebibliography}

\end{document}
